
\chapter{Rozwinięcie}

Rozdziały dokumentujące pracę własną studenta: opisujące ideę, sposób lub metodę 
rozwiązania postawionego problemu oraz rozdziały opisujące techniczną stronę rozwiązania 
--- dokumentacja techniczna, przeprowadzone testy, badania i uzyskane wyniki. 

Praca musi zawierać elementy pracy własnej autora adekwatne do jego wiedzy praktycznej uzyskanej w
okresie studiów. Za pracę własną autora można uznać np.: stworzenie aplikacji informatycznej lub jej
fragmentu, zaproponowanie algorytmu rozwiązania problemu szczegółowego, przedstawienie projektu 
np.~systemu informatycznego lub sieci komputerowej, analizę i ocenę nowych technologii lub rozwiązań
informatycznych wykorzystywanych w przedsiębiorstwach, itp. 

Autor powinien zadbać o właściwą dokumentację pracy własnej obejmującą specyfikację założeń i 
sposób realizacji poszczególnych zadań
wraz z ich oceną i opisem napotkanych problemów. W przypadku prac o charakterze 
projektowo-implementacyjnym, ta część pracy jest zastępowana dokumentacją techniczną i użytkową systemu. 

W pracy \textbf{nie należy zamieszczać całego kodu źródłowego} opracowanych programów. Kod źródłowy napisanych
programów, wszelkie oprogramowanie wytworzone i wykorzystane w pracy, wyniki przeprowadzonych
eksperymentów powinny być umieszczone np. na płycie CD, stanowiącej dodatek do pracy.

\section*{Styl tekstu}

Należy\footnote{Uwagi o stylu pochodzą częściowo ze stron prof. Macieja Drozdowskiego~\cite{Drozdowski2006}.} 
stosować formę bezosobową, tj.~\emph{w pracy rozważono ......, 
w ramach pracy zaprojektowano ....}, a nie: \emph{w pracy rozważyłem, w ramach pracy zaprojektowałem}. 
Odwołania do wcześniejszych fragmentów tekstu powinny mieć następującą postać: ,,Jak wspomniano wcześniej, ....'', 
,,Jak wykazano powyżej ....''. Należy unikać długich zdań. 

Niedopuszczalne są zwroty używane w języku potocznym. W pracy należy używać terminologii informatycznej, która ma 
sprecyzowaną treść i znaczenie. 

Niedopuszczalne jest pisanie pracy metodą \emph{cut\&paste}, bo jest to plagiat i dowód intelektualnej indolencji autora.
Dane zagadnienie należy opisać własnymi słowami. Zawsze trzeba powołać się na zewnętrzne źródła. 

\section{Wkład Jakuba Górniaka --- szkic}

Roboczy zakres własnych prac do późniejszego rozwinięcia. Wyniki pomiarów
dotyczą opisanych poniżej prób, a nie końcowej oceny modeli na niezależnym
zbiorze testowym.

\begin{itemize}
	\item \textbf{Przygotowanie zbioru danych:} zebranie i ręczne oznaczenie
	ok. 2700 linii na 194 stronach notatek anatomicznych i dokumentów medycznych.
	Przygotowanie transkrypcji i ramek pisma odręcznego, również na formularzach
	zawierających druk. Charakterystyka trudności: tekst wokół rysunków,
	kratka, nieregularne i nachodzące linie oraz polskie znaki.

	\item \textbf{Stworzenie narzędzia do labelowania:} autorska aplikacja
	webowa do oznaczania ramek i transkrypcji, oparta na Flask i HTML5 Canvas.
	Obsługa obrazów i wielostronicowych PDF, obracanie stron z przeliczaniem
	współrzędnych, edycja ramek, zoom, przesuwanie, cofanie i skróty klawiszowe.
	Zapis stron, wycinków i etykiet oraz ponowne wczytywanie anotacji.

	\item \textbf{Automatyczne wspomaganie anotacji:} propozycje ramek
	i tekstu do ręcznej weryfikacji, przetwarzanie dokumentów w tle,
	monitorowanie postępu i anulowanie oraz ponowny odczyt po poprawieniu
	ramki. Wymienne moduły detekcji i OCR, zbiorcze zatwierdzanie propozycji
	oraz rejestrowanie pochodzenia anotacji jako ręcznych, automatycznych,
	zatwierdzonych lub poprawionych.

	\item \textbf{Wersjonowanie i udostępnianie danych:} migracja obrazów
	z Gita do DVC, początkowo z Google Drive, następnie DagsHub. Scalanie
	rozproszonych partii z zachowaniem istniejących anotacji i kopiami
	bezpieczeństwa. Automatyzacja publikowania skryptem oraz obejście
	problemu dużych obrazów w historii Gita przez oddzielne publikowanie
	wskaźnika danych.

	\item \textbf{Porównanie detektorów:} implementacja heurystyki
	morfologicznej z usuwaniem kratki, scalaniem fragmentów i ustalaniem
	kolejności czytania oraz porównanie jej z CRAFT, Surya i YOLO.
	Przygotowanie benchmarku, strojenie progów CRAFT i sprawdzenie ustawień
	na innych stronach. Porównanie metryk pokrycia i dopasowania IoU 1:1.

	\item \textbf{Wyniki detekcji:} wspólny test na 30 stronach i 571 liniach
	referencyjnych, z ziarnem losowania 0 i wyłączeniem kart pacjenta.
	Przy progu pokrycia 0.5 CRAFT z domyślnymi progami uzyskał F1 0.722,
	Surya 0.656, YOLO 0.649, a heurystyka morfologiczna 0.204.
	Przy dopasowaniu 1:1 z progiem IoU 0.5 najlepszy wynik uzyskał YOLO
	(F1 0.296 wobec 0.195 dla domyślnego CRAFT).
	Orientacyjny czas wynosił 0.6 s na stronę dla YOLO i 3.7 s dla CRAFT.
	Analiza wpływu metryki na ranking i kompromisu między jakością a czasem.

	\item \textbf{Eksperymenty i dostrajanie OCR:} sprawdzenie Qwen3-VL
	przez Ollama i odrzucenie tej konfiguracji z powodu pustych odpowiedzi
	i opóźnienia. Przygotowanie eksportu danych treningowych i skryptu
	dostrajania TrOCR z monitorowaniem CER i wczesnym zatrzymaniem.
	Dwa warianty treningu uzyskały walidacyjny CER 0.112 i 0.047;
	różniły się bazą językową oraz ustawieniami treningu, dlatego poprawy
	nie można przypisać wyłącznie polskiej bazie modelu.

	\item \textbf{Analiza błędów kadrowania:} porównanie rozpoznawania
	wycinków z ramek ręcznych i automatycznych oraz zależności CER
	od dopasowania ramki. Zaobserwowanie znacznie lepszych wyników dla
	dobrze dopasowanych wycinków i wprowadzenie augmentacji kadrowania,
	aby przybliżyć trening do warunków użycia.

	\item \textbf{Integracja i sprawdzenie modeli zespołu:} podłączenie
	dostarczonego YOLO do wykrywania linii oraz Surya z adapterem LoRA
	do odczytu wycinków, bez przypisywania sobie treningu tych modeli.
	Sprawdzenie rozmiaru wejścia, progu pewności, tłumienia nakładających
	się detekcji i preprocessingu. F1 pokrycia YOLO wzrosło z 0.304
	dla 1024 px bez preprocessingu do 0.649 dla 512 px z preprocessingiem.
	Na 40 losowych wycinkach średni CER liczony po liniach spadł z 0.237
	dla bazowej Surya do 0.175 z adapterem. Sprawdzenie całego przepływu:
	detekcja, wycinanie linii i OCR. Przygotowanie dokumentacji rozwiązania
	i eksperymentów do późniejszego opisu w pracy.
\end{itemize}

